\documentclass{article}
\usepackage[utf8]{inputenc}
\usepackage[T1]{fontenc}
\usepackage{geometry}
\usepackage{graphicx}
\usepackage{float}
\usepackage{hyperref}
\usepackage[english,indonesian]{babel}

\title{Optimasi Pemrosesan Template LaTeX Menggunakan Pendekatan Berbasis AI}
\author{Academic Support Agent}
\date{\today}

\begin{document}
\maketitle

\begin{abstract}

\end{abstract}

\chapter{PENDAHULUAN}


\chapter{TINJAUAN PUSTAKA}


% Original content (commented out):
% \documentclass[11pt]{article} % Can change font size here

% ------------------------ CAN DELETE TO DECLUTTER ------------------------
% CHANGE COMPILER TO XELATEX - IMPORTANT OTHERWISE ERROR!
% This is because the package for Fancy headers requires an alternative compiler
% Check the correct compiler is chosen by clicking 'Menu' in the top left, under settings select % the drop down box next to 'compiler' and select XeLaTex.
% https://www.overleaf.com/learn/latex/Tutorials - ALL TUTORIALS FOR LATEX 
% Any problems that are not contained within the link above, please consult coding forums such as % StackOverflow OR email 2237553@swansea.ac.uk
% ------------------------------------------------------------------------

\usepackage{graphicx} % Required for inserting images
\usepackage{biblatex}
\usepackage{geometry}
\geometry{a4paper,
        total={170mm,245mm},
        left=20mm,
        top=20mm,
        }
\usepackage{fancyhdr}
\pagestyle{fancy}
\fancyhead{}
\fancyhead[L]{\noindent{\Large \textbf{College of Engineering}}\\
                {\normalsize \textbf{Coursework Submission Sheet}}}
\fancyhead[R]{\includegraphics{Images/swansealogo.png}}
\renewcommand{\headrulewidth}{0pt}
\setlength{\headsep}{2cm}
\fancyheadoffset{1cm}
\usepackage{fontspec}
\setmainfont{Arial} % Change to whatever your brief states you should use
\usepackage{biblatex}
% \bibliographystyle{} Change referncing style here (currently Vancouver (numbered)) 
\addbibresource{references.bib}

\begin{document}

\pagestyle{fancy}
\input{titlepage}

% input pages here
% create a file (e.g. 'intro.text' - can place in Writing folder) and call it here using...
% \input{Writing/****.tex}

%\printbibliography - remove comment to see references (after using \cite{})

\end{document}



\end{document}